% Part 6 – Implementation of Noise Conditional Score Networks

% ---------------------------------------------------------------------------
% In this section we translate the theoretical insights from Part 5 into a
% concrete implementation of noise conditional score networks (NCSNs) for
% the MNIST dataset.  We describe the data preprocessing steps, noise
% schedule, neural network architecture, loss function, training loop and
% sampling algorithm.  We also discuss how to extend the model to
% condition on class labels.  Throughout, we emphasise reproducibility and
% clarity, providing code snippets where appropriate.

\section{Noise--Conditional Score Networks: Implementation}

\subsection{Data Preparation and Noise Ladder}

As with the energy–based model, we begin by loading the MNIST dataset.
However, to stabilise score estimation we normalise the pixel values to
the range $[-1,1]$ using the transformation $\tilde{\mathbf{x}} = 2\mathbf{x}
- 1$, where $\mathbf{x}$ is the original image scaled to $[0,1]$.  We
employ the same \texttt{DataLoader} configuration as before but
optionally drop the labels if training the unconditional model.  For the
conditional model we retain the digit labels $y\in\{0,\dots,9\}$.

A key component of NCSNs is the \emph{noise ladder}.  We choose a
geometric sequence of $L=10$ noise levels defined by
\begin{equation}
    \sigma_i = \sigma_{\max} \left(\frac{\sigma_{\min}}{\sigma_{\max}}\right)^{i/L},
    \quad i=1,\dots,L,
    \label{eq:sigma-schedule}
\end{equation}
with $\sigma_{\max}=30.0$ and $\sigma_{\min}=0.01$.  Sampling uniformly
from this ladder ensures that each scale is used equally often during
training.  We cache the noise levels on the GPU for efficiency.

\subsection{Architecture: U--Net with Adaptive ResBlocks and FiLM}

The score network must map a noisy image $\tilde{\mathbf{x}}$ and a noise
level $\sigma$ to a tensor of the same shape as the input, representing
the estimated score.  Following the assignment, we implement a U–Net
style architecture with skip connections and noise conditioning via
feature–wise linear modulation (FiLM).  The network comprises three
encoder blocks, a bottleneck, and two decoder blocks.  Each block is an
\emph{AdaptiveResBlock} that consists of a convolution, group
normalisation, an activation, and FiLM modulation.

\paragraph{Gaussian Fourier embeddings.}  To embed the noise level
$\sigma$ into a high–dimensional vector, we apply a Gaussian Fourier
projection.  Specifically, we sample a fixed random matrix $\mathbf{B}
\in\mathbb{R}^{m\times 1}$ from a Gaussian distribution and compute
\begin{equation}
    \phi(\sigma) = \big[ \cos(\mathbf{B}\sigma),\, \sin(\mathbf{B}\sigma)\big],
\end{equation}
yielding a $2m$–dimensional embedding.  A small multilayer perceptron
(MLP) then projects $\phi(\sigma)$ to a 256–dimensional vector.  This
embedding captures periodic structure across different noise levels and
serves as the conditioning signal for FiLM.

\paragraph{FiLM modulation.}  Given a feature map $\mathbf{h}\in
\mathbb{R}^{C\times H\times W}$ at a particular layer and the projected
noise embedding $\mathbf{e}$, we compute scale and shift vectors
$\gamma,\beta\in\mathbb{R}^C$ using a linear transformation of
$\mathbf{e}$ and apply
\begin{equation}
    \text{FiLM}(\mathbf{h},\gamma,\beta) = \gamma \odot \mathbf{h} + \beta,
\end{equation}
where $\odot$ denotes element–wise multiplication.  In the
class–conditional model we further add a class embedding to the noise
embedding before computing $\gamma$ and $\beta$.

\paragraph{Encoder and decoder.}  The network begins with a $3\times 3$
convolution that maps the single–channel input to 64 channels.  Two
encoder blocks follow, each downsampling the spatial resolution by a
factor of two via average pooling.  At the bottleneck, the channel
dimension is 256.  The decoder mirrors the encoder: each decoder block
upsamples via nearest–neighbour interpolation, concatenates the
corresponding encoder features (skip connection), and applies an
AdaptiveResBlock to reduce the channels.  A final $3\times 3$ convolution
with a single output channel produces the score estimate.

\paragraph{Group normalisation and SiLU.}  Within each block we use
group normalisation with 32 groups and the SiLU activation function
(\texttt{torch.nn.SiLU}) rather than BatchNorm or ReLU.  Group norm is
preferable when batch sizes are small because it normalises across
channels instead of examples.  The SiLU activation is smooth and
differentiable everywhere, which aids gradient–based training.

\paragraph{Class conditioning.}  To extend the network to generate
digits of a specific class, we introduce an \texttt{nn.Embedding} layer
that maps each label $y\in\{0,\dots,9\}$ to a 256–dimensional vector.
This embedding is added to the noise embedding before FiLM projection in
each block.  During training we provide the correct label alongside the
image and noise level; during sampling we select the desired label and
freeze it throughout the annealing process.

\subsection{Loss Function: Weighted Denoising Score Matching}

The training objective follows the weighted DSM formulation described in
Part~5.  Given a mini–batch of images, we randomly select a noise level
$\sigma_i$ from the ladder and corrupt the images according to
Eq.~\eqref{eq:sigma-schedule}.  We then compute the score network's output
and compare it to the target $-\boldsymbol{\epsilon}/\sigma_i$.  The
weighted DSM loss is
\begin{equation}
    \mathcal{L}(\theta) = \frac{1}{2} \sigma_i^2 \Big\| \mathbf{s}_\theta(\tilde{\mathbf{x}},\sigma_i) + \frac{\boldsymbol{\epsilon}}{\sigma_i} \Big\|_2^2.
\end{equation}
In code this becomes:
\begin{lstlisting}[language=Python]
def weighted_dsm_loss(model, x, sigmas):
    batch_size = x.size(0)
    # Sample a noise level per example
    idx = torch.randint(0, len(sigmas), (batch_size,), device=x.device)
    sigma = sigmas[idx].view(-1, 1, 1, 1)
    # Corrupt the input
    noise = torch.randn_like(x)
    x_noisy = x + sigma * noise
    # Predict the score
    score = model(x_noisy, sigma)
    target = -noise / sigma
    loss = 0.5 * (sigma.squeeze() ** 2) * (score - target).pow(2).reshape(batch_size, -1).mean(dim=1)
    return loss.mean()
\end{lstlisting}
We optimise this loss using the Adam optimiser with a learning rate of
$2\times 10^{-4}$ and train for 30 epochs.  We use a batch size of 64
and shuffle the training data each epoch.  The learning rate, number of
noise levels, and number of epochs can be adjusted; our choices balance
training time and sample quality.

\subsection{Training Loop and Logging}

The training loop for NCSN parallels that of the EBM but with notable
differences:

\begin{enumerate}[leftmargin=*,label=\arabic*.]
  \item \textbf{Noise sampling.}  For each mini–batch we draw a noise
    level index and create noisy inputs using Eq.~\eqref{eq:dsm-perturb}.
  \item \textbf{Forward pass.}  We pass the noisy images and the noise
    level (and optionally the label) through the network to obtain the
    predicted scores.
  \item \textbf{Loss computation.}  We compute the weighted DSM loss
    described above.  Because the model outputs a tensor of shape
    $(B,1,28,28)$, we reshape it into vectors and average over pixels.
  \item \textbf{Optimiser step.}  We call \texttt{loss.backward()} and
    update the network parameters using Adam.
  \item \textbf{Logging and sample generation.}  After each epoch we
    record the average loss and perform annealed Langevin dynamics to
    generate a batch of samples from pure noise.  These samples allow us
    to monitor the evolution of sample quality during training.  We
    normalise the generated images back to $[0,1]$ for visualisation.
\end{enumerate}

In the conditional setting the labels are provided alongside each image
and used to look up a class embedding.  The rest of the training loop
remains unchanged.

\subsection{Annealed Langevin Dynamics Sampler}

Once the network has been trained, we generate samples by running ALD
starting from a Gaussian distribution with variance $\sigma_L^2$.  The
procedure is summarised in the following pseudocode:
\begin{lstlisting}[language=Python]
def annealed_langevin_sampler(model, sigmas, shape, y=None, steps_per_level=150):
    # Initialise from a high-noise Gaussian
    x = torch.randn(*shape, device=sigmas.device) * sigmas[-1]
    for sigma in reversed(sigmas):
        step_size = 2e-5 * (sigma / sigmas[0])**2
        for _ in range(steps_per_level):
            grad = model(x, torch.tensor([sigma], device=x.device), y)
            noise = torch.randn_like(x)
            x = x + step_size * grad + (2 * step_size)**0.5 * noise
        x = x.clamp(-1.0, 1.0)
    return x
\end{lstlisting}
Here \texttt{shape} specifies the batch size and image dimensions
(e.g., \texttt{(16,1,28,28)}).  The step size is scaled by $(\sigma/\sigma_1)^2$
to ensure stability across noise levels.  The function returns images in
the normalised space $[-1,1]$; we map them back to $[0,1]$ when saving.

\subsection{Conditional Generation and Label Embeddings}

To enable class–conditional generation we augment the score network with
an embedding layer \texttt{nn.Embedding(10,256)}.  At each noise level
the label embedding is added to the projected noise embedding before the
FiLM transformation.  During training we sample labels uniformly and
train the network to predict the score conditioned on the true label.  At
inference time we set $y$ to the desired digit class and run ALD as
before.  By varying the random seed one obtains diverse samples within
the specified class.  In Part~7 we visualise the result by producing a
$10\times10$ grid in which each row corresponds to a digit class.

\subsection{Hyperparameters and Practical Tips}

Table~\ref{tab:hyperparams-ncsn} summarises the hyperparameters used in
our implementation.  As with the EBM, hyperparameters influence sample
quality and training stability.  Increasing the number of noise levels or
steps per level improves fidelity at the cost of runtime.  Larger batch
sizes stabilise the loss but require more memory.

\begin{table}[H]
    \centering
    \caption{Key hyperparameters for the MNIST NCSN.}
    \label{tab:hyperparams-ncsn}
    \begin{tabular}{lll}
        \toprule
        Parameter & Value & Description \\
        \midrule
        $\sigma_{\max}$ & 30.0 & Largest noise level \\
        $\sigma_{\min}$ & 0.01 & Smallest noise level \\
        Number of levels $L$ & 10 & Noise scales in ladder \\
        Epochs & 30 & Training duration \\
        Batch size & 64 & Images per mini–batch \\
        Learning rate & $2\times 10^{-4}$ & Adam step size \\
        Steps per level & 150 & ALD iterations per noise level \\
        Step size scaling $\eta_0$ & $2\times 10^{-5}$ & Base step size for ALD \\
        Embedding dimension & 256 & Size of noise and label embeddings \\
        Conditional & Yes/No & If labels are used \\
        \bottomrule
    \end{tabular}
\end{table}

\subsection{Implementation Summary}

In this part we described how to implement a noise–conditional score
network for MNIST.  We normalised the data to $[-1,1]$, defined a
geometric noise ladder, designed a U–Net architecture with FiLM
conditioning, and employed weighted denoising score matching as the loss.
We presented a training loop that samples noise levels and updates the
network via Adam, and provided a sampling algorithm based on annealed
Langevin dynamics.  A class–conditional extension was implemented by
adding a label embedding.  These details prepare the ground for the
results and analysis presented in the next part.